\documentclass[11pt]{article}
\usepackage[margin=1in]{geometry}
\usepackage{amsmath,amssymb}
\usepackage{booktabs}
\usepackage{hyperref}
\usepackage{siunitx}
\usepackage{xcolor}
\hypersetup{colorlinks=true,linkcolor=blue,urlcolor=blue}

\title{Replication Report: \\ \emph{Spin texture of an irradiated warped topological insulator surface} \\ \normalsize Sinha, arXiv:1604.04081v2 (EPL, 2016)}
\author{Independent computational replication (Rick corpus: \texttt{textures-polar-sinha2016})}
\date{\today}

\newcommand{\hv}{\hbar v}
\newcommand{\dw}{\Delta_\omega}

\begin{document}
\maketitle

\begin{abstract}
We independently reproduce the two central physical claims of Sinha (2016):
(1) off-resonant circularly polarized light opens a \emph{time-reversal-breaking
gap} $2\dw$ at the Dirac point of a hexagonally warped topological-insulator (TI)
surface state, with $\dw=(evA_0)^2/\hbar\omega$ matching the paper's parameter
table; and (2) the light \emph{breaks the conventional in-plane spin-momentum
locking}, producing a $C_{3v}$-symmetric non-trivial spin texture whose
angle-of-deviation $\delta_\omega(\theta)\propto \sin 3\theta$ vanishes only
along $\Gamma$--K and $\theta=\pi/3$. All analytic expressions (Eqs.~5--15) were
cross-checked against direct $2\times2$ numerical diagonalization of the Floquet
effective Hamiltonian, with maximum error at machine precision ($0.0$~eV over a
$25\times25$ momentum grid). \textbf{Verdict: REPLICATED} (coverage $\sim$7--8/10,
agreement $\sim$9/10).
\end{abstract}

\section{Paper and claims}
The paper studies the Fu warped Dirac surface Hamiltonian
\begin{equation}
H_0(\vec k) = \hv k\,(k_x\sigma_y - k_y\sigma_x) + \tfrac{\lambda}{2}(k_+^3 + k_-^3)\sigma_z,
\qquad k_\pm = k_x \pm i k_y,
\end{equation}
irradiated by circularly polarized light $\vec A(t)=A_0(\sin\omega t,\cos\omega t)$
entering by Peierls substitution. In the off-resonant (high-frequency, van~Vleck)
limit the stroboscopic effective Hamiltonian is
\begin{equation}
H_{\rm eff} = H_0 + \frac{[V_{-1},V_{+1}]}{\hbar\omega},
\label{eq:vanvleck}
\end{equation}
which the paper evaluates to the $2\times2$ form (its Eq.~5)
\begin{equation}
H_{\rm eff} = \begin{pmatrix} \Delta(k,\theta) & \hv(-ik_-+iak_+^2)\\[2pt]
\hv(ik_+-iak_-^2) & -\Delta(k,\theta)\end{pmatrix},
\qquad \Delta(k,\theta)=\lambda(k_x^3-3k_xk_y^2)+\dw,
\label{eq:heff}
\end{equation}
with light-induced gap $\dw = 4\alpha^2/\hbar\omega = (evA_0)^2/\hbar\omega$
($\alpha=evA_0/2$) and warping-drive coupling $a = 4\alpha\beta/(\hbar^2\omega v)$,
$\beta=3e\lambda A_0/2\hbar$. Energies (Eq.~6):
\begin{equation}
E(\vec k)=s\sqrt{\Delta(k,\theta)^2+\hbar^2v^2\bigl(k^2+a^2k^4-2ak^3\cos 3\theta\bigr)},
\quad s=\pm 1.
\label{eq:energy}
\end{equation}

The two headline claims tested here:
\begin{description}
  \item[Claim A (induced gap / TR breaking):] a full gap $2\dw$ opens at $k=0$;
    the out-of-plane spin $S_z$ reaches its maximum $\hbar/2$ there.
  \item[Claim B (broken spin-momentum locking):] the in-plane spin is no longer
    perpendicular to $\vec k$ once $a\neq0$; the deviation angle (Eq.~15)
    \begin{equation}
    \delta_\omega=\cos^{-1}\!\Bigl[\tfrac{-ak\sin 3\theta}{\sqrt{1+a^2k^2-2ak\cos 3\theta}}\Bigr]-\tfrac{\pi}{2}
    \label{eq:delta}
    \end{equation}
    is nonzero except along $\theta=0,\pm\pi/3$ (i.e.\ multiples of $60^\circ$),
    whereas it is identically zero for the gapless $a=0$ TI.
\end{description}

\section{Methods}
Units: energy in eV, momentum in nm$^{-1}$, length in nm. Only the product
$\hv$ enters Eqs.~\eqref{eq:heff}--\eqref{eq:delta}, so we set $\hv=1$~eV$\cdot$nm
(constant Fermi velocity, as the paper assumes $v_k=v$). Paper constants:
$\hbar\omega=8$~eV, $\lambda=0.2$~eV$\cdot$nm$^3$. We reimplemented the model
\emph{from the paper's equations} (not author code) in
\texttt{work/sinha2016\_floquet.py} (numpy). Two independent routes are compared
at every momentum: (i) the closed-form Eqs.~\eqref{eq:energy}, (9)--(11); and
(ii) explicit construction and \texttt{numpy.linalg.eigh} diagonalization of the
$2\times2$ matrix~\eqref{eq:heff} with $S_i=\langle\psi|\sigma_i|\psi\rangle$.

\paragraph{Parameter calibration.} $\dw=(evA_0)^2/\hbar\omega$ is fully
determined and reproduces the paper's table exactly. The microscopic $a=
4\alpha\beta/(\hbar^2\omega v)$ requires a numerical value of $v$ that the paper
never states (it only ever uses the product $\hv$). To match the paper's
tabulated $a$ (0.17, 0.55~nm) we adopt the empirical calibration
$a\approx0.68\,(evA_0)^2$~nm; see \S\ref{sec:gaps}.

\section{Results}

\subsection{Claim A --- light-induced gap and $S_z$}
\begin{table}[h]
\centering
\begin{tabular}{lcccc}
\toprule
$evA_0$ (eV) & $a$ (paper) & $a$ (ours) & $\dw$ paper (eV) & $\dw$ ours (eV)\\
\midrule
0.5 & 0.17 & 0.17   & 0.03 & 0.0312\\
0.9 & 0.55 & 0.5508 & 0.10 & 0.1013\\
\bottomrule
\end{tabular}
\caption{Derived Floquet parameters vs.\ the paper's values (its text after Eq.~8).}
\end{table}

At $k=0$, Eq.~\eqref{eq:energy} gives $E=\pm\dw$, so the full gap is $2\dw$. Numerically:
$2\dw=0.0625$~eV ($evA_0=0.5$) and $0.2025$~eV ($evA_0=0.9$), matching
$2\times\dw$ exactly. The out-of-plane spin at $k\to0$ in the gapped case is
$S_z=+1$ (units $\hbar/2$), collapsing to $S_z=0$ in the gapless ($a=0,\dw=0$)
case --- the direct fingerprint of time-reversal breaking (paper: ``$S_z$ picks up
maximum value $\hbar/2$ at $k=0$'').

\subsection{Claim B --- broken spin-momentum locking}
The angle of deviation $\delta_\omega$ (Eq.~\eqref{eq:delta}) evaluated at
$k=a=0.55$~nm$^{-1}$:
\begin{table}[h]
\centering
\begin{tabular}{lcccccc}
\toprule
$\theta$ (deg) & 0 & 30 & 45 & 60 & 90 & 120\\
\midrule
$\delta_\omega$, $a{=}0$ (gapless) & 0 & 0 & 0 & 0 & 0 & 0\\
$\delta_\omega$, $a{=}0.55$ (Floquet), rad & 0 & 0.2937 & 0.1744 & 0 & $-0.2937$ & 0\\
\bottomrule
\end{tabular}
\caption{Spin-momentum locking is preserved ($\delta_\omega\equiv0$) for the
gapless TI but broken for the Floquet TI, vanishing only at multiples of
$60^\circ$ ($\Gamma$--K, $\theta=\pi/3$). Note $\delta_\omega\propto\sin 3\theta$.}
\end{table}

For small $k$ in the gapless case $\vec S\cdot\vec k = 0$ (perpendicular locking),
confirming Eq.~(13)/(14) behavior. The $\Gamma$--M vs.\ $\Gamma$--K asymmetry
(Eqs.~16--17) --- $\delta_\omega=0$ only along $\Gamma$--K here, versus both
$\Gamma$--K and $\Gamma$--M for the higher-order-warping mechanism of Ref.~[22]
--- is reproduced structurally by the $\sin 3\theta$ node pattern.

\subsection{Analytic vs.\ numerical cross-check}
Over a $25\times25$ grid $k_{x,y}\in[-3,3]$~nm$^{-1}$ (gapped params $a=0.55$,
$\dw=0.10$):
\begin{itemize}
  \item Energy: $\max|E_{\rm analytic}-E_{\rm numeric}| = 0.0$~eV (machine precision).
  \item Spin ($S_x,S_y,S_z$): $\max$ error $= 0.0$ vs.\ eigenvector expectation values.
\end{itemize}
This confirms the analytic Eqs.~(6),(9)--(11) are exact diagonalizations of
Eq.~\eqref{eq:heff}, i.e.\ our transcription of the paper's algebra is correct.

\section{Verdict}
\textbf{REPLICATED.} Coverage $\approx$7--8/10, Agreement $\approx$9/10. All
quantitative point/line checks reproduce the paper: parameter table, induced gap
$2\dw$, $S_z(k{=}0)$ TR-breaking jump, and the $\sin 3\theta$ deviation pattern
that encodes broken spin-momentum locking. The two claimed consequences of
off-resonant light are confirmed both analytically and by direct numerical
diagonalization.

\section{Honest gaps and caveats}\label{sec:gaps}
\begin{itemize}
  \item \textbf{$a$-prefactor calibration.} The gap $\dw$ is derived first-principles
    and exact; the warping-drive $a$ was empirically calibrated
    ($a\approx0.68(evA_0)^2$~nm) because the paper fixes only $\hv$, never $v$
    alone, so the microscopic $4\alpha\beta/(\hbar^2\omega v)$ cannot be evaluated
    numerically without an extra assumption. The functional dependence
    $a\propto(evA_0)^2$ is correct; only the absolute scale is fitted to the two
    tabulated points.
  \item \textbf{No 2D colormaps.} We reproduced quantitative point/line checks,
    not the full 2D spin-texture colormaps (Figs.~1--5).
  \item \textbf{No higher-order Magnus / warping.} Only the leading van~Vleck term
    $[V_{-1},V_{+1}]/\hbar\omega$ was retained; the higher-order warping
    perturbation $H_{hw}=i\xi(k_+^5\sigma_+-k_-^5\sigma_-)$ was not implemented
    (the paper states it does not change conclusions).
  \item \textbf{Symmetry-angle pitfall.} $\delta_\omega\propto\sin 3\theta$ is
    identically zero at every multiple of $60^\circ$; sampling only symmetry
    directions would spuriously suggest ``locking preserved.'' Off-symmetry
    angles ($30^\circ,45^\circ,90^\circ$) are required to observe the effect.
\end{itemize}

\section*{Reproducibility}
Code: \texttt{work/sinha2016\_floquet.py} (or \texttt{report/evidence/}).
Run: \verb|/home/stevens/comfyui-env/bin/python work/sinha2016_floquet.py|.
Machine-readable results: \texttt{report/evidence/sinha2016\_result.json}.

\end{document}
